\section{The introspective $\neexp$ protocol}\label{sec:big-neexp-protocol}

In this question, we give the complete short-question, introspective $\neexp$ protocol.
The goal is a protocol for $\succinctsquared$ instances of size~$s_{\mathrm{inst}}$
with $\poly(s_{\mathrm{inst}})$ question length and running time and $\poly(2^{s_{\mathrm{inst}}})$ answer length and running time.
Our construction will be an introspective version of the classical PCP construction from \Cref{sec:classical-pcp},
in which we replace the low-degree tests and simultaneous low-degree tests
with our introspective low-degree test and introspective simultaneous low-degree test.

We summarize the protocol here.
Given the $\succinctsquared$ instance~$\calC_{\mathrm{inst}}$
of size~$s_{\mathrm{inst}}$,
let $\calC$ be the size-$s$, $(3n+3)$-input $\succinct$ instance it succinctly represents,
where~$s$ and~$n$ are roughly exponential in~$s_{\mathrm{inst}}$.
Following \Cref{sec:classical-pcp},
we would like the introspective prover to sample strings $\bx_1, \bx_2, \bx_3 \in \F_q^m$ and $(\bb, \bw) \in \F_q^{3+s}$,
which they should return to the verifier.
In addition, they should return the evaluations $g(\bx_1), g(\bx_2), g(\bx_3)$ and $c_1(\bx, \bb, \bw), \ldots, c_{m'}(\bx, \bb, \bw)$,
where $g$ and the $c_i$'s are purported degree-$d$ polynomials.
This suggests using the following registers:
\begin{equation*}
\ket{\mathrm{EPR}_q^m}_{\reg{1}}
\otimes \ket{\mathrm{EPR}_q^m}_{\reg{2}}
\otimes \ket{\mathrm{EPR}_q^m}_{\reg{3}}
\otimes \ket{\mathrm{EPR}_q^{3+s}}_{\reg{4}}.
\end{equation*}
The difficulty in this protocol is ensuring that the polynomials involved are low-degree.
To begin, we can run the introspective low-degree test on the first register,
which guarantees that $g(\bx_1)$ corresponds to a low-degree polynomial. 
Doing the same on registers~$2$ and~$3$ would guarantee
the functions evaluated on~$\bx_2$ and~$\bx_3$ are also low-degree polynomials,
but it would not guarantee that the prover is using \emph{same} low-degree polynomial~$g$ on all three.
Instead, we run the introspective intersecting lines test twice,
ensuring that prover evaluates~$\bx_1$, $\bx_2$, and~$\bx_3$ using the same function~$g$.

Next, we consider the coefficient polynomials $c_1, \ldots, c_{m'}$.
They are evaluated on the concatenated outputs of the four registers, i.e.\ the string $(\bx, \bb, \bw)$.
As a result, we view the four registers as a single superregister of length $m' = 3 m + 3 + s$,
and we would like to perform the introspective simultaneous low-degree test on this superregister.
However, this test requires two additional superregisters of length $m'$ to serve as the direction registers.
As a result, the shared state between the two provers will be of the following form:
\begin{multline*}
\ket{\psi} = 
(\ket{\mathrm{EPR}_q^m}_{\reg{1}})
\otimes \ket{\mathrm{EPR}_q^m}_{\reg{2}}
\otimes \ket{\mathrm{EPR}_q^m}_{\reg{3}}
\otimes \ket{\mathrm{EPR}_q^{3+s}}_{\reg{4}})_{\reg{Super Reg 1}}\\
\otimes (\ket{\mathrm{EPR}_q^{m'}}_{\reg{5}})_{\reg{Super Reg 2}}
\otimes (\ket{\mathrm{EPR}_q^{m'}}_{\reg{6}})_{\reg{Super Reg 3}}
\otimes \ket{\mathrm{aux}}_{\reg{aux}}.
\end{multline*}

Having checked that the provers' functions are low-degree,
we conclude with a consistency check between~$g$ and the~$c_i$'s
to ensure that they encode a satisfying assignment to our $\succinct$.
In \Cref{sec:classical-pcp},
this was done by the ``formula test",
i.e.\ the check that $\mathrm{sat}_{\psi, g}(\bx, \bb, \bw) = \mathrm{zero}_{H, c}(\bx, \bb, \bw)$.
Here, this will be accomplished by an introspective version of this test,
in which the provers sample~$\bx$, $\bb$, and $\bw$ themselves.
Passing this test with high probability proves
that $\calC_{\mathrm{inst}}$ is a YES instance of the $\succinctsquared$ problem.

This section is organized as follows.
In \Cref{sec:computing-registers}, we will discuss the register parameters algorithm,
needed for the register compiler from \Cref{sec:register-overview}.
Next, \Cref{sec:introspective-formula} introduces the introspective formula game.
Finally, \Cref{sec:the-whole-banana} completes the construction and gives the introspective $\neexp$ game.

\subsection{Computing the register parameters}\label{sec:computing-registers}

Given the $\succinctsquared$ instance~$\calC_{\mathrm{inst}}$
of size~$s_{\mathrm{inst}}$,
let $\calC$ be the size-$s$, $(3n+3)$-input $\mathsf{Succinct}$-$\mathsf{3Sat}$ instance it succinctly represents.
To compile our protocol to one sound against general provers,
we need a register parameters algorithm which runs in time $\poly(s_{\mathrm{inst}})$ (\Cref{def:reg-params-generator}).
As described above, the register parameters will be simple functions of the numbers~$s$ and~$n$
(for example $m$, a simple function of~$n$ to be determined later).
However, $s$ and~$n$ themselves may not be easy to compute,
as the natural way of computing them involves first computing~$\calC$,
a time $2^{s_{\mathrm{inst}}}$ task.
We solve this by ``guessing" values for these numbers which are guaranteed to be larger than the actual values,
and then later ``fixing" the circuit~$\calC$ so that it actually has the guessed input length and size.
This is detailed in the following definition.

\begin{definition}\label{def:register-parameters}
Let $\calC_{\mathrm{inst}}$ be a size-$s_{\mathrm{inst}}$ instance of the $\succinctsquared$ problem.
\begin{enumerate}
\item Let $\calC$ be the size-$s$ $\mathsf{Succinct}$-$\mathsf{3Sat}$ instance it succinctly represents.
This circuit takes inputs $i, j, k$, each of some length~$n$, and bits $b_1, b_2, b_3$.
Then $s$ and $n$ can both be trivially upper-bounded by $N := 2^{s_{\mathrm{inst}}}$.
\item Consider a new circuit $\calC_{\mathrm{pad}}$ with inputs $i, j, k \in \{0, 1\}^{N}$ and $b \in \{0, 1\}^3$. We write $i = (i_1, i_2)$, where $i_1$ is of length $N - n$ and $i_2$ is of length $n$, and likewise for~$j$ and~$k$. Let this circuit act as follows:
\begin{itemize}
\item[$\circ$] Compute the $\lor$ of the bits in $i_1$, $j_1$, and~$k_1$. Output~$0$ if this is~$1$.
\item[$\circ$] Otherwise, output $\calC_{\mathrm{dec}}(i_2, j_2, k_2, b_1, b_2, b_3)$.
\end{itemize}
As defined, this circuit has size $s + 3 (N - n) + 2 \leq 4 N =: S$,
and we will pad it with additional gates in a trivial manner so that it has exactly~$S$ gates.
It can be checked that it succinctly represents the same $\sat$ formula as $\calC_{\mathrm{dec}}$.
\end{enumerate}
We set $\mathrm{PadC}(\calC_{\mathrm{inst}}) := \calC_{\mathrm{pad}}$, $\mathrm{PadN}(\calC_{\mathrm{inst}}) := N$,
and $\mathrm{PadS}(\calC_{\mathrm{inst}}):= 4 \cdot N$.
We note that given $\calC_{\mathrm{inst}}$, the value of~$N$ is efficiently computable.
\end{definition}

\subsection{An introspective formula game}\label{sec:introspective-formula}

In this section, we introduce the ``introspective formula game".
This game is the introspective version of the formula check in \Cref{sec:classical-pcp},
in which we check $\mathrm{sat}_{\psi, g}(\bx, \bb, \bw) = \mathrm{zero}_{H, c}(\bx, \bb, \bw)$
on a randomly chosen point $(\bx, \bb, \bw)$ in $\F_q^{m'}$.
Prior to stating the introspective formula game,
we will begin by recalling what this notation means.

Let $\calC_{\mathrm{inst}}$ be a size-$(s_{\mathrm{inst}})$ $\succinctsquared$ instance.
Let $\calC = \mathrm{PadC}(\calC_{\mathrm{inst}})$ be a $\succinct$ instance,
and let $n = \mathrm{PadN}(\calC_{\mathrm{inst}})$ and $s = \mathrm{PadS}(\calC_{\mathrm{inst}})$.
Then $\calC$ is a size-$s$, $(3n+3)$-variable circuit
which is a YES instance of the $\succinct$ problem if and only if $\calC_{\mathrm{inst}}$
is a YES instance of the $\succinctsquared$  problem.
Introduce $h = 2^{t_1}$, $q = 2^{t_2}$, and $m$ such that $N=2^n$, $h$, $q$, and $m$ are exactly admissible parameters (\Cref{def:admissible}).
Set $n' = n + 3 + s$ and $m' = m + 3 + s$.
We also recall the following pieces of notation.
\begin{itemize}
\item[$\circ$] (\Cref{def:canonical-low-degree}): Write $H := H_{t_1, t_2}$.
\item[$\circ$]  (\Cref{def:formula-function}): Given a function $g:\F_q^m \rightarrow \F_q$,
recall the notation $\mathrm{sat}_{\psi,g} := \mathrm{sat}_{\psi,g, n, t_1, t_2}$.
\item[$\circ$]  (\Cref{prop:coefficient-polys}): Writing $H_{\mathrm{zero}} = H^{3m} \otimes \{0, 1\}^{3 + s}$.
Given $c_1, \ldots, c_{m'} : \F_q^{m'} \rightarrow \F_q$, recall the notation $\mathrm{zero}_{H,c} = \mathrm{zero}_{H_{\mathrm{zero}}, c}$.
\end{itemize}

Before stating the introspective formula game, we must first dispense with the following annoying technicality.

\begin{notation}\label{not:annoying-technicality}
In the classical case (\Cref{sec:classical-pcp}),
we have a fixed proof which contains fixed functions which may or may not be low-degree.
In the quantum case, however, we are dealing not with a fixed proof but an interactive prover,
and the formula prover may not respond based on fixed functions (their responses might be randomized, for example).
To account for this, we modify the definitions of $\mathrm{sat}$ and $\mathrm{zero}$ as follows.
First, we recall the notation $g_\psi := g_{\psi, n, t_1, t_2}$ (\Cref{def:encoded-function}). 
\begin{itemize}
	\item[$\circ$] Given $\nu_1, \nu_2, \nu_3 \in \F_q$, define
		\begin{equation*}
			\mathrm{sat}_{\psi, \nu}(x, b, w) := g_{\psi}(x, b, w) \cdot (\nu_1 - b_1) (\nu_2 - b_2) (\nu_3 - b_3).
		\end{equation*}
	\item[$\circ$] Given $\mu_1, \ldots, \mu_{m'} \in \F_q$, define
		\begin{equation*}
			\mathrm{zero}_{H, \mu}(x) = \sum_{i=1}^{m'} \mathrm{zero}_{(H_{\mathrm{zero}})_i}(x_i) \cdot \mu_i,
		\end{equation*}
		where by definition $(H_{\mathrm{zero}})_i = H$ for $i \in [3m]$ and $(H_{\mathrm{zero}})_i = \{0, 1\}$ otherwise.
\end{itemize}
We note that if there is a function~$g$ such that $\nu_i = g(x_i)$, then $\mathrm{sat}_{\psi, \nu} = \mathrm{sat}_{\psi, g}$.
Similarly, if there are functions $c_1, \ldots, c_{m'}$ such that $\mu_i = c_i(x)$, then $\mathrm{zero}_{H, \mu} = \mathrm{zero}_{H, c}$.
\end{notation}

Now we state the introspective formula game.

\begin{definition}\label{def:formula-game}
Let $\calC_{\mathrm{inst}}$ be a size-$(s_{\mathrm{inst}})$ $\succinctsquared$ instance.
Let $n = \mathrm{PadN}(\calC_{\mathrm{inst}})$ and $s = \mathrm{PadS}(\calC_{\mathrm{inst}})$.
Suppose $n$, $h = 2^{t_1}$, $q = 2^{t_2}$, and $m$ are exactly admissible parameters.
The \emph{introspective formula game},
denoted $\game := \game_{\mathrm{IntroForm}}(\calC_{\mathrm{inst}}, h, q, m)$, is defined in \Cref{fig:formula-game}.
This is a $\lambda_{\calC_{\mathrm{inst}}, q} := (4, \ell, q)$-register game, for $\ell = (m, m, m, 3 + s)$.
Furthermore,
\begin{equation*}
\qlength{\game} = O(1), \quad
\alength{\game} = O(m' \log(q)),
\end{equation*}
\begin{equation*}
\qtime{\game} = O(1), \quad
\atime{\game} = \poly(s, n, n', h, q, m').
\end{equation*}
\end{definition}
{
\floatstyle{boxed} 
\restylefloat{figure}
\begin{figure}
Flip an unbiased coin $\bb \sim \{0, 1\}$.
	\begin{itemize}
	\item[$\circ$] Player~$\bb$: Give ($Z$, $Z$, $Z$, $Z$, ``formula"); 
				receive $\bu_1, \bu_2, \bu_3, (\bb, \bw)$ and $\bnu_1, \bnu_2, \bnu_3$ and $\bmu_1, \ldots, \bmu_{M'}$.
	\end{itemize}
	Compute~$\mathrm{sat}_{\psi, \bnu}(\bu, \bb, \bw)$ and~$\mathrm{zero}_{H, \bmu}(\bu, \bb, \bw)$.
	Accept if they are equal.
	\caption{The game $\game_{\mathrm{IntroForm}}(\calC_{\mathrm{inst}}, h, q, m)$.\label{fig:formula-game}}
\end{figure}
}

\begin{notation}
In the case when a prover is given the question $(Z, Z, Z, Z, \mathrm{``formula"})$, we refer to it as the \emph{formula prover}. 
It has the following intended behavior.
\begin{enumerate}
\setcounter{enumi}{2}
\item \textbf{Formula prover:}
	\begin{description}[align=left]
	\item [Input:] Pauli basis queries $(Z, Z, Z, Z)$ and auxiliary query ``formula".
	\item [Output:] Strings $u_1, u_2, u_3 \in\F_q^{m}$ and $(b, w) \in \F_q^{3+s}$.
			Three numbers $\nu_1,\nu_2, \nu_3 \in \F_q$
			and $m'$ numbers $\mu_1, \ldots, \mu_{m'} \in \F_q$.
	\item [Goal:] The prover should act as follows.
		\begin{itemize}
		\item[$\circ$] The prover sets $\nu_1 = g(u_1)$, $\nu_2 = g(u_2)$, $\nu_3 = g(u_3)$,
					where $g:\F_q^m \rightarrow \F_q$ is a global degree-$d_1$ polynomial selected independently of~$u$.
		\item[$\circ$] They then set $\mu_i = c_i(u_1, u_2, u_3, b, w)$,
					where for each $i$,
					$c_i:\F_q^{m'} \rightarrow \F_q$ is a global degree-$d_2$ polynomial selected independently of $(u, b, w)$.
		\end{itemize}
	\end{description}
\end{enumerate}
Here, $d_1$ and~$d_2$ are polynomial degrees which will be selected later.
We will also refer to the \emph{formula prover's measurement}, 
which refers to the measurement $\{F_{u, b, w, \nu_i,\mu_j}\}$ such that
\begin{equation*}
F_{u, b, w, \nu_1, \nu_2, \nu_3, \mu_1, \ldots, \mu_{m'}}
= M^{Z, Z, Z, Z, \mathrm{``formula"}}_{u, b, w, \nu_1, \nu_2, \nu_3, \mu_1, \ldots, \mu_{m'}}.
\end{equation*}
\end{notation}


We begin by showing the completeness case of the introspective formula game.

\begin{proposition}[Introspective formula game completeness]\label{prop:formula-game-completeness}
Suppose~$\calC_{\mathrm{inst}}$ is a YES instance of the $\succinctsquared$ problem.
Let $a:\{0,1\}^n \rightarrow \{0, 1\}$ be a satisfying assignment to the $\mathsf{3Sat}$ instance it encodes,
and let $g:=g_a:\F_q^m \rightarrow \F_q$ be its low-degree encoding.
Let $c_1, \ldots, c_{m'}:\F_q^{m'} \rightarrow \F_q$ be the coefficient polynomials guaranteed to make 
$\mathrm{sat}_{\psi, g} = \mathrm{zero}_{H_{\mathrm{zero}}, c}$ by \Cref{prop:coefficient-polys}.
Both $g$ and the $c_i$'s are degree-$O(hn')$ polynomials.
Consider the $\lambda_{\calC_{\mathrm{inst}},q}$-register strategy $(\psi, A)$ with no auxiliary register in which
\begin{equation*}
A_{u, b, w, \nu, \mu} =  \tau^Z_{u_1} \otimes \tau^Z_{u_2} \otimes \tau^Z_{u_3} \otimes \tau^Z_{b, w}  \cdot \bone[\nu_i = g(u_i), \mu_j = c_j(u, b, w)],
\end{equation*}
where the indices range over $i \in [3]$ and $j \in [m']$.
Then this strategy passes $\game_{\mathrm{IntroForm}}(\calC_{\mathrm{inst}}, h, q, m)$ with probability~$1$.
\end{proposition}
\begin{proof}
This game is simply the oracularized version of the formula check in the classical PCP.
The proposition follows from the discussion in \Cref{sec:the-pcp}.
\end{proof}

Our next lemma covers the soundness case of the introspective formula game.
It concerns provers of a particular form,
namely those whose measurements correspond to low-degree polynomials.
We show that if there exists such a prover with nonnegligible value, then the formula must be satisfiable.

\begin{lemma}[Formula game partial soundness]\label{lem:polynomial-means-good}
Let $\calC_{\mathrm{inst}}$ be a $\succinctsquared$ instance, and set $\game := \game_{\mathrm{IntroForm}}(\calC_{\mathrm{inst}}, h, q, m)$.
Let $\calS = (\psi, A)$ be a $\lambda_{\calC_{\mathrm{inst}}, q}$-register strategy.
Consider a measurement on the auxiliary register
\begin{equation*}
G = \{G_{g, c_1, \ldots, c_{m'}}\}
\end{equation*}
with outcomes degree-$d_1$ polynomials $g:\F_q^m \rightarrow \F_q$
and degree-$d_2$ polynomials $c_1, \ldots, c_{m'}: \F_Q^{m'} \rightarrow \F_q$.
Suppose~$A$ has the following form: for each $u$, $b$, $w$, $\nu$, and $\mu$,
\begin{equation}\label{eq:weird-form}
A_{u, b, w, \nu, \mu} = \tau^Z_{u_1} \otimes \tau^Z_{u_2} \otimes \tau^Z_{u_3} \otimes \tau^Z_{b, w} \otimes \left(G_{[g(u_i)= \nu_i,  c_j(u, b, w)= \mu_j]}\right)_{\reg{aux}},
\end{equation}
where the subscript of the~$G$ measurement ranges over all $i \in [3]$ and $j \in [m']$.
If the probability $\calS$ passes $\calG$ is at least
\begin{equation*}
\frac{\max\{O(h n') + 3 d_1, h + d_2\}}{q},
\end{equation*}
then $\psi$ is satisfiable.
\end{lemma}
\begin{proof}
Consider the following three step strategy:
\begin{enumerate}
\item Measure the auxiliary register with $\{G_{g, c_1, \ldots, c_{m'}}\}$, receiving functions $\bg$, $\bc_1$, \ldots, $\bc_{m'}$.
\item Measure the EPR registers in the~$Z$ basis, receiving $\bu$, $\bb$, and~$\bw$.
\item Output $\bu$, $\bb$, $\bw$, $\bg(\bu_1)$, $\bg(\bu_2)$, $\bg(\bu_3)$ and $\bc_1(\bu,\bb, \bw)$ through $\bc_{M'}(\bu, \bb, \bw)$.
\end{enumerate}
This passes the formula game with probability $\valstrat{\calG}{\calS}$.
Then there exists functions $g$, $c_1$, \ldots, $c_{m'}$ such that conditioned on measuring them in step one,
this strategy passes with probability at least $\valstrat{\calG}{\calS}$.
By the remark at the end of \Cref{not:annoying-technicality},
this is the probability that
\begin{equation*}
\mathrm{sat}_{\psi, g}(\bx, \bb, \bw) = \mathrm{zero}_{H, c}(\bx, \bb, \bw),
\end{equation*}
where $(\bx, \bb, \bw)$ is drawn from $\F_q^{m'}$ uniformly at random.
The lemma follows from \Cref{lem:the-grand-finale}.
\end{proof}


\subsection{The complete introspective protocol}\label{sec:the-whole-banana}

In this section, we introduce the introspective protocol for $\neexp$ and prove its correctness.
The introspective $\neexp$ protocol builds on top of the introspective formula game
by using a series of introspective low-degree tests to ensure that the formula prover satisfies the condition in \Cref{eq:weird-form}.
Having done this, we can then apply \Cref{lem:polynomial-means-good},
ensuring that if a strategy passes with high probability, then the instance is satisfiable.

\begin{definition}\label{def:the-protocol-to-end-all-protocols}
Let $\calC_{\mathrm{inst}}$ be a size-$(s_{\mathrm{inst}})$ $\succinctsquared$ instance.
Let $n = \mathrm{PadN}(\calC_{\mathrm{inst}})$ and $s = \mathrm{PadS}(\calC_{\mathrm{inst}})$.
The verifier chooses $h = 2^{t_1}$, $q = 2^{t_2}$, $m$, and $d$
such that $m$, $h$, $q$, and $m$ are exactly admissible parameters satisfying
\begin{equation*}
h = \Theta(n),
\quad
m = \Theta\left(\frac{n}{\log(n)}\right),
\quad
q = \poly(n),
\quad
d = O(hn') = O(n^2).
\end{equation*}
(We will choose the polynomial for~$q$ in \Cref{thm:main-result-of-this-part} below.)
The verifier sets $\lambda = (6, \ell, q)$,
where $\ell = (m,m,m, 3 + s, m', m')$.

We begin by instantiating the following list of subroutines.
\begin{itemize}
\item[$\circ$]
Let $\lambda_{\mathrm{LD}} = (3, m, q)$ be register parameters.
Let $\game_{\mathrm{LD}}$ be a copy of $\game_{\mathrm{IntroLowDeg}}(\lambda_{\mathrm{LD}}, d)$,
	using register~$1$ as the point register and registers~$2$ and~$3$ as the direction registers.
	Write $\mathsf{Points}_1$ for the points prover.
\item[$\circ$]
Let $\lambda_{\mathrm{IL}} = (2, m, q)$ be register parameters.
Let $\game_{\mathrm{IL1}}$ be a copy of $\game_{\mathrm{IntroIntersect}}(\lambda_{\mathrm{IL}}, d)$ on registers~$1$ and~$2$
	whose points prover for register~$1$ is $\mathsf{Points}_1$ from $\game_{\mathrm{LD}}$.
	Write $\mathsf{Points}_2$ for the points prover on register~$2$.
\item[$\circ$] Let $\game_{\mathrm{IL2}}$ be a copy of
	$\game_{\mathrm{IntroIntersect}}(\lambda_{\mathrm{IL}}, d)$ on registers~$1$ and~$3$
	whose points prover for register~$1$ is $\mathsf{Points}_1$ from $\game_{\mathrm{LD}}$.
	Write $\mathsf{Points}_3$ for the points prover on register~$3$.
\item[$\circ$] Let $\game_{\mathrm{F}}$ be a copy of $\game_{\mathrm{IntroForm}}(\calC_{\mathrm{inst}}, h, q, m)$ on registers~$1$, $2$, $3$, and~$4$.
	Write $\mathsf{Formula}$ for the formula prover.
\item[$\circ$]
Let $\lambda_{\mathrm{LDSUP}} = (3, m', q)$ be register parameters.
Let $\game_{\mathrm{LDSUP}}$ be a copy of $\game_{\mathrm{IntroLowDeg}}(\lambda_{\mathrm{LDSUP}}, d, 3+m')$,
applied to the following three superregisters:
registers~$1$ through~$4$ are combined into the point superregister,
register~$5$ is used as the first direction superregister,
and register~$6$ is used as the second direction superregister.
In addition, use $\mathsf{Formula}$ from $\game_{\mathrm{form}}$ as its points prover.
\end{itemize}
Then the \emph{introspective $\neexp$ game}, denoted $\game_{\mathrm{Intro}\neexp}(\calC_{\mathrm{inst}})$, is defined in \Cref{fig:big-neexp}.
\end{definition}

{
\floatstyle{boxed} 
\restylefloat{figure}
\begin{figure}
With probability~$\tfrac{1}{9}$ each, perform one of the following nine tests.
\begin{enumerate}
	\item \textbf{Low degree test:} Play $\game_{\mathrm{LD}}$.
	\item \textbf{Intersecting lines test 1:} Play $\game_{\mathrm{IL1}}$.
	\item \textbf{Intersecting lines test 2:} Play $\game_{\mathrm{IL2}}$.
	\item \textbf{Simultaneous low degree test:} Play $\game_{\mathrm{LDSUP}}$.
	\item \textbf{Formula test:} Player $\game_{\mathrm{F}}$.
\end{enumerate}
For the remaining tests, flip an unbiased coin $\bb \sim \{0, 1\}$.
Assign the first role to Player~$\bb$ and the second role to Player~$\overline{\bb}$.
\begin{enumerate}
\setcounter{enumi}{5}
	\item \textbf{Consistency test 1:} 
		\begin{itemize}
		\item[$\circ$] $\mathsf{Points}_1$: Receive $\bnu$.
		\item[$\circ$] $\mathsf{Formula}$: Receive $\bnu_1$.
		\end{itemize}
		Accept if $\bnu = \bnu_1$.
	\item \textbf{Consistency test 2:} 
		\begin{itemize}
		\item[$\circ$] $\mathsf{Points}_2$: Receive $\bnu$.
		\item[$\circ$] $\mathsf{Formula}$: Receive $\bnu_2$.
		\end{itemize}
		Accept if $\bnu = \bnu_2$.
	\item \textbf{Consistency test 3:} 
		\begin{itemize}
		\item[$\circ$] $\mathsf{Points}_3$: Receive $\bnu$.
		\item[$\circ$] $\mathsf{Formula}$: Receive $\bnu_3$.
		\end{itemize}
		Accept if $\bnu = \bnu_3$.
	\item \textbf{Consistency test 4:} 
		\begin{itemize}
		\item[$\circ$] $\mathsf{Formula}$: Receive $\bnu_1, \bnu_2, \bnu_3$ and $\bmu_1, \ldots, \bmu_{m'}$.
		\item[$\circ$] $\mathsf{Formula}$: Receive $\bnu_1', \bnu_2', \bnu_3'$ and $\bmu_1', \ldots, \bmu_{m'}'$.
		\end{itemize}
		Accept if $\bnu_i = \bnu_i'$ and $\bmu_j = \bmu_j'$ for all $i \in [3]$, $j \in [m']$.
	\end{enumerate}
	\caption{The game $\game_{\mathrm{Intro}\neexp}(\calC_{\mathrm{inst}})$.\label{fig:big-neexp}}
\end{figure}
}

The main result \Cref{part:neexp} is the following theorem.

\begin{theorem}\label{thm:main-result-of-this-part}
Let $\calC_{\mathrm{inst}}$ be a size-$(s_{\mathrm{inst}})$ $\succinctsquared$ instance.
Let $q$ be a sufficiently large $\poly(n)$ and $\eps > 0$ a sufficiently small constant such that \Cref{eq:formula-success-probability} is at least $\tfrac{1}{2}$
and \Cref{eq:other-expression-dont-even-know-what-to-call-this} is less than $\tfrac{1}{2}$.
Write $\game := \game_{\mathrm{Intro}\neexp}(\calC_{\mathrm{inst}})$.
\begin{itemize}
\item[$\circ$] \textbf{Completeness:}
		Suppose $\calC_{\mathrm{inst}}$ encodes a satisfiable formula.
		Then there is a value-$1$ $\lambda$-register strategy for $\game$ with no auxiliary register.
\item[$\circ$] \textbf{Soundness:}
		If there is a $\lambda$-register strategy for $\game$ with value at least $1-\eps$,
		then $\calC_{\mathrm{inst}}$ encodes a satisfiable formula.
\end{itemize}
Furthermore,
\begin{equation*}
\qlength{\game} = O(1), \quad
\alength{\game} = \poly(2^{s_\mathrm{inst}}),
\end{equation*}
\begin{equation*}
\qtime{\game} = O(1), \quad
\atime{\game} = \poly(2^{s_\mathrm{inst}}).
\end{equation*}
\end{theorem}

\begin{proof}
The question lengths and question runtimes are both $O(1)$ because all involved subtests have $O(1)$ question complexity.
The answer lengths and question runtimes are both $\poly(2^{s_{\mathrm{inst}}})$
because all our parameters are at most polynomial in $n = 2^{s_{\mathrm{inst}}}$,
and the question lengths and question runtimes of all involved subtests are polynomial in these parameters.

We name the measurements used by the provers as follows.
\begin{equation*}
\mathsf{Points}_1: A,\quad
\mathsf{Points}_2: B,\quad
\mathsf{Points}_3: C,\quad
\mathsf{Formula}: F.
\end{equation*}
We will write the identity matrix on registers~$5$ and~$6$ as $I_{\reg{5, 6}} := I_{\reg{5}} \otimes I_{\reg{6}}$.

\paragraph{Completeness.}
Suppose $\calC_{\mathrm{inst}}$ encodes a satisfiable formula.
By \Cref{prop:formula-game-completeness}, 
there are degree-$d$ polynomials $g:\F_q^m \rightarrow \F_q$ and $c_1, \ldots, c_{m'}:\F_q^{m'} \rightarrow \F_q$ such that if we define
\begin{equation*}
F_{u, b, w, \nu, \mu} =  \tau^Z_{u_1} \otimes \tau^Z_{u_2} \otimes \tau^Z_{u_3} \otimes \tau^Z_{b, w} \otimes I_{\reg{5,6}} \cdot \bone[\nu_i = g(u_i), \mu_j = c_j(u, b, w)],
\end{equation*}
then this strategy passes the formula test with probability~$1$.
We extend this strategy to the remaining measurements as follows.
\begin{equation*}
A_{u_1, \nu_1}
= \tau^Z_{u_1} \otimes I_{\reg{2}} \otimes I_{\reg{3}} \otimes I_{\reg{4}}\otimes I_{\reg{5,6}}\cdot \bone[\nu_1 = g(u_1)],
\end{equation*}
\begin{equation*}
B_{u_2, \nu_2}
= I_{\reg{1}} \otimes \tau^Z_{u_2} \otimes I_{\reg{3}} \otimes I_{\reg{4}} \otimes I_{\reg{5,6}} \cdot\bone[\nu_2 = g(u_2)],
\end{equation*}
\begin{equation*}
C_{u_3, \nu_3}
= I_{\reg{1}} \otimes I_{\reg{2}} \otimes \tau^Z_{u_3} \otimes I_{\reg{4}} \otimes I_{\reg{5, 6}}\cdot \bone[\nu_3 = g(u_3)].
\end{equation*}
By the completeness of the introspective low-degree and intersecting lines tests, these can be extended to a strategy which passes the whole test with probability~$1$.

\paragraph{Soundness.}
Throughout this proof, we use $\delta(\eps)$ to represent functions of the form
\begin{equation*}
\delta(\eps) = \poly(\eps, m \cdot d/q^e),
\end{equation*}
where $e > 0$ is an absolute constant.

\paragraph{Low-degree tests.}
The strategy passes the introspective low-degree test with probability $1-\delta(\eps)$.
Applying~\Cref{thm:big-planes-low-degree}, there is a measurement $G = \{G_g\}$ in $\mathrm{PolyMeas}(m, d, q)$ such that
\begin{equation*}
(A_{u_1, \nu_1})_{\reg{Alice}} \otimes I_{\reg{Bob}}
\approx_{\delta(\eps)}
\left(\tau^Z_{u_1} \otimes I_{\reg{2}} \otimes I_{\reg{3}} \otimes I_{\reg{4}} \otimes I_{\reg{5,6}} \otimes (G_{[g(u_1) = \nu_1]})_{\reg{aux}}\right)_{\reg{Alice}} \otimes I_{\reg{Bob}}.
\end{equation*}
By \Cref{fact:approx-delta-game-value}, we can assume this holds with equality
with a loss of only $\delta(\eps)$ in the game value.
In addition, by \Cref{thm:naimark}, we can assume that the~$G$ measurements are all projective.

Next, the strategy passes the two introspective intersecting lines tests with probability $1-\delta(\eps)$ each.
By~\Cref{thm:big-transfer}, this implies that
\begin{equation}\label{eq:whatevs}
(B_{u_2, \nu_2})_{\reg{Alice}} \otimes I_{\reg{Bob}}
\approx_{\delta(\eps)}
\left(I_{\reg{1}} \otimes \tau^Z_{u_2} \otimes I_{\reg{3}} \otimes I_{\reg{4}} \otimes I_{\reg{5,6}} \otimes (G_{[g(u_2) = \nu_2]})_{\reg{aux}}\right)_{\reg{Alice}} \otimes I_{\reg{Bob}},
\end{equation}
\begin{equation}\label{eq:whatevs-dos}
(C_{u_3, \nu_3})_{\reg{Alice}} \otimes I_{\reg{Bob}}
\approx_{\delta(\eps)}
\left(I_{\reg{1}} \otimes I_{\reg{2}} \otimes \tau^Z_{u_3} \otimes I_{\reg{4}} \otimes I_{\reg{5,6}} \otimes (G_{[g(u_3) = \nu_3]})_{\reg{aux}}\right)_{\reg{Alice}} \otimes I_{\reg{Bob}}.
\end{equation}

Similarly, the strategy passes the introspective simultaneous low-degree test with probability $1-\delta(\eps)$.
Applying~\Cref{thm:big-simultaneous-planes-low-degree}, there is a measurement $J = \{J_{f_1, f_2, f_3, c_1, \ldots, c_{m'}}\}$ in $\mathrm{PolyMeas}(m',d, q, 3+m')$ such that
\begin{multline}\label{eq:whatevs-tres}
(F_{u, b, w, \nu, \mu})_{\reg{Alice}} \otimes I_{\reg{Bob}}\\
\approx_{\delta(\eps)}
\left(\tau^Z_{u_1} \otimes \tau^Z_{u_2} \otimes \tau^Z_{u_3} \otimes \tau^Z_{b, w}
	\otimes I_{\reg{5,6}} \otimes (J_{[f_i(u, b, w) = \nu_i, c_j(u, b, w) = \mu_j]})_{\reg{aux}}\right)_{\reg{Alice}} \otimes I_{\reg{Bob}},
\end{multline}
where the subscript of the~$J$ measurement ranges over all $i \in [3]$ and $j \in [m']$.
By \Cref{fact:approx-delta-game-value}, we can assume \Cref{eq:whatevs,eq:whatevs-dos,eq:whatevs-tres} holds with equality
with a loss of only $\delta(\eps)$ in the game value.
In addition, by \Cref{thm:naimark}, we can assume that the~$J$ measurements are all projective.

\paragraph{Consistency tests.}
The strategy passes the four consistency tests with probability $1-\delta(\eps)$ each, implying
\begin{equation*}
(F_{u_1, \nu_1})_{\reg{Alice}} \otimes I_{\reg{Bob}} \simeq_{\delta(\eps)} I_{\reg{Alice}} \otimes (A_{u_1, \nu_1})_{\reg{Bob}},
\end{equation*}
\begin{equation*}
(F_{u_2, \nu_2})_{\reg{Alice}} \otimes I_{\reg{Bob}} \simeq_{\delta(\eps)} I_{\reg{Alice}} \otimes (B_{u_2, \nu_2})_{\reg{Bob}},
\end{equation*}
\begin{equation*}
(F_{u_3, \nu_3})_{\reg{Alice}} \otimes I_{\reg{Bob}} \simeq_{\delta(\eps)} I_{\reg{Alice}} \otimes (C_{u_3, \nu_3})_{\reg{Bob}},
\end{equation*}
\begin{equation*}
(F_{u, b, w, \nu, \mu})_{\reg{Alice}} \otimes I_{\reg{Bob}} \simeq_{\delta(\eps)} I_{\reg{Alice}} \otimes (F_{u, b, w, \nu, \mu})_{\reg{Bob}}.
\end{equation*}
By introspection~(\Cref{fact:introspective-consistency}), these imply the following statements:
\begin{equation*}
(J_{[f_1(u, b, w) = \nu_1]})_{\reg{Alice}} \otimes I_{\reg{Bob}} \simeq_{\delta(\eps)} I_{\reg{Alice}} \otimes (G_{[g(u_1) = \nu_1]}),
\end{equation*}
\begin{equation*}
(J_{[f_2(u, b, w) = \nu_2]})_{\reg{Alice}} \otimes I_{\reg{Bob}} \simeq_{\delta(\eps)} I_{\reg{Alice}} \otimes (G_{[g(u_2) = \nu_2]}),
\end{equation*}
\begin{equation*}
(J_{[f_3(u, b, w) = \nu_3]})_{\reg{Alice}} \otimes I_{\reg{Bob}} \simeq_{\delta(\eps)} I_{\reg{Alice}} \otimes (G_{[g(u_3) = \nu_3]}),
\end{equation*}
\begin{equation*}
(J_{[c_j(u, b, w) = \mu_j]})_{\reg{Alice}} \otimes I_{\reg{Bob}} \simeq_{\delta(\eps)} I_{\reg{Alice}} \otimes (J_{[c_j(u, b, w) = \mu_j]}),
\end{equation*}
where the subscript of the~$J$ measurement ranges over all $i \in [3]$ and $j \in [m']$.
Here, these statements are with respect to the strategy's auxiliary state and to the uniform distribution on $(\bu, \bb, \bw) \in \F_q^{m'}$.

Now we apply \Cref{fact:low-degree-sandwich}.
To do so, let us specify the sets $\calG_i$ and the distance parameter.
The three sets $\calG_2$, $\calG_3$, and $\calG_4$ will just contain all degree-$d$ polynomials $g:\F_q^{m'} \rightarrow \F_q$.
(Note that we can view the outputs of $G_g$ as degree-$d$ polynomials which disregard all of their input $(u, b, w)$ aside from one of the three strings $u_1$, $u_2$, or $u_3$.)
By Schwarz-Zippel, these have distance at least $1-d/q$.
The remaining set, $\calG_1$, is defined as follows:
for each tuple of degree-$d$ polynomials $c_1, \ldots, c_{m'}$,
it contains a function~$c$ defined as $c(u, b, w) = (c_1(u, b, w), \ldots, c_{m'}(u, b, w))$.
Any two nonequal $c, c' \in \calG_1$ have some coordinate~$i$ in which $c_i \neq c_i'$,
and on this coordinate alone they will have distance at least $1-d/q$ by Schwarz-Zippel.
Thus, $c$ and $c'$ have distance at least $1-d/q$.


Define the measurement $\{K_{g, c_1, \ldots, c_{m'}}\}$ as 
\begin{equation*}
K_{g, c_1, \ldots, c_{m'}} := G_g \cdot J_{c_1, \ldots, c_{m'}} \cdot G_g,
\end{equation*}
Then \Cref{fact:low-degree-sandwich} implies that
\begin{equation*}
(J_{[f_i(u, b, w) = \nu_i, c_j(u, b, w) = \mu_j]})_{\reg{Alice}} \otimes I_{\reg{Bob}} \simeq_{\delta(\eps)} I_{\reg{Alice}} \otimes (K_{[g(u_i) = \nu_i, c_j'(u, b, w) = \mu_j]}),
\end{equation*}
where the subscripts range over all $i \in [3]$ and $j \in [m']$.
By introspection~(\Cref{fact:introspective-consistency}), this implies that
\begin{multline}\label{eq:cant-believe-it-but-we-did-it}
(F_{u, b, w, \nu, \mu})_{\reg{Alice}} \otimes I_{\reg{Bob}}\\
\simeq_{\delta(\eps)}
\left(\tau^Z_{u_1} \otimes \tau^Z_{u_2} \otimes \tau^Z_{u_3} \otimes \tau^Z_{b, w}
	\otimes I_{\reg{5,6}} \otimes (K_{[g(u_i) = \nu_i, c_j(u, b, w) = \mu_j]})_{\reg{aux}}\right)_{\reg{Alice}} \otimes I_{\reg{Bob}},
\end{multline}
where the subscripts range over all $i \in [3]$ and $j \in [m']$.
By \Cref{fact:approx-delta-game-value}, we can assume \Cref{eq:cant-believe-it-but-we-did-it} holds with equality
with a loss of only $\delta(\eps)$ in the game value.

\paragraph{Formula test:}
At this point, the formula prover's strategy~$F$ satisfies the condition in \Cref{eq:weird-form} with $d_1 = d_2 = d$.
In addition, it passes the introspective formula test with probability
\begin{equation}\label{eq:formula-success-probability}
1-\poly(\eps, m \cdot d/q),
\end{equation}
which by our setting of parameters is at least $\tfrac{1}{2}$.
Finally, our setting of parameters also implies that
\begin{equation}\label{eq:other-expression-dont-even-know-what-to-call-this}
\frac{\max\{O(h n') + 3d, h + d\}}{q}
\end{equation}
is less than $\tfrac{1}{2}$.
As a result, we can apply \Cref{lem:polynomial-means-good} to conclude that~$\psi$ is satisfiable.
\end{proof}

\Cref{thm:main-result-of-this-part} only proves soundness of the introspective $\neexp$ protocol against $\lambda$-register strategies.
Our last step is to compile this protocol into one which is sound against \emph{general} strategies,
while only slightly increasing the question length.

\begin{corollary}\label{cor:succinct-sat-protocol-with-big-answer-size}
There is an absolute constant $\eps > 0$ such that the following is true.
Let $\calC_{\mathrm{inst}}$ be a size-$(s_{\mathrm{inst}})$ $\succinctsquared$ instance.
Then there exists a game $\game := \game_{\mathrm{Intro}\neexp}(\calC_{\mathrm{inst}})$ with the following properties.
\begin{itemize}
\item[$\circ$] \textbf{Completeness:}
		Suppose $\calC_{\mathrm{inst}}$ encodes a satisfiable formula.
		Then there is a value-$1$ real commuting EPR strategy for $\game$.
\item[$\circ$] \textbf{Soundness:}
		If there is a strategy for $\game$ with value at least $1-\eps$,
		then $\calC_{\mathrm{inst}}$ encodes a satisfiable formula.
\end{itemize}
Furthermore,
\begin{equation*}
\qlength{\game} = O(s_{\mathrm{inst}}), \quad
\alength{\game} = \poly(2^{s_{\mathrm{Inst}}}),
\end{equation*}
\begin{equation*}
\qtime{\game} = O(s_{\mathrm{inst}}), \quad
\atime{\game} = \poly(2^{s_{\mathrm{inst}}}).
\end{equation*}
\end{corollary}
\begin{proof}
Let $n = \mathrm{PadN}(\calC_{\mathrm{inst}})$ and $s = \mathrm{PadS}(\calC_{\mathrm{inst}})$.
Set $m = \Theta(n/\log(n))$ and $q = \poly(n)$,
as in \Cref{def:the-protocol-to-end-all-protocols}.
Set $\lambda = (6, \ell, q)$,
where $\ell = (m, m, m, 3+s, 3m + 3 + s, 3m + 3 + s)$.
Then $\game_{\mathrm{Intro}\neexp}(\calC_{\mathrm{inst}})$
is a $\lambda$-register game.
Furthermore, by \Cref{def:register-parameters},
the register parameters are computable in time $\poly(s_{\mathrm{inst}})$.

Let $\eps >0$ be as in \Cref{thm:main-result-of-this-part}, and select a constant $\eps' > 0$ and $\eta_1, \ldots, \eta_{6} = 1/\poly(n)$
such that $\delta(\eps') \leq \eps$, where $\delta(\eps') = \poly(\eps', \eta_1, \ldots, \eta_{6})$ is as in \Cref{thm:uniform-registers}.
Now, if we apply \Cref{thm:uniform-registers}, it gives us a game~$\game$ with the following properties.
\begin{itemize}
\item[$\circ$] If $\calC$ is a ``Yes" instance, then $\game_{\mathrm{Intro}\neexp}(\calC_{\mathrm{inst}})$ has a value-$1$ strategy with no auxiliary state, which implies that $\game$ has a value-$1$ commuting EPR strategy.
\item[$\circ$] If $\calC$ is a ``No" instance, then every $\lambda$-register strategy for
		$\game_{\mathrm{Intro}\neexp}(\calC_{\mathrm{inst}})$ has value less than $1-\eps$.
		By our choice of parameters, this is less than $1-\delta(\eps')$, which implies that $\game$ has no strategy with value $1-\eps'$.
\end{itemize}
Furthermore, $\log(m) = O(\log(n)) = O(s_{\mathrm{inst}})$ and $\log(s) = O(\log(n)) = O(s_{\mathrm{inst}})$, giving us our desired question complexities,
and $\poly(m) = \poly(n) = \poly(2^{s_{\mathrm{inst}}})$ and $\poly(s) = \poly(n) = \poly(2^{s_\mathrm{inst}})$, giving us our desired answer complexities.
\end{proof}

%%% Local Variables:
%%% mode: latex
%%% TeX-master: "main"
%%% End:
%%%---------- close: big-neexp-protocol

\part{Answer reduction}

\label{part:answer}

%%%%%%%%%%% jump to error-correcting-code-test
%%%---------- open: error-correcting-code-test
%!TEX root = main.tex

